Based on the above, highly cited papers are considered to exhibit the following characteristics:
(1) Highly cited papers tend to be cited across diverse themes. This is suggested by the most frequent themes in Table \ref{tab:topmesh}, and the degree of thematic transition illustrated in Figure \ref{fig:meshpivot} further supports this finding.
(2) Highly cited papers have focused research themes, and the papers that cite them also exhibit concentrated thematic profiles. This is indicated by Figure \ref{fig:arr2}.
(3) The thematic diversity of highly cited papers is low from a synchronic perspective, but not necessarily so from a diachronic one. This is because, while Figure \ref{fig:arr2} shows lower entropy compared to other groups, Figure \ref{fig:gr15} reveals no clear distinction when compared with those groups.

